\hypertarget{ux4f7fux7528ux5143ux7c7b}{%
\subsection{使用元类}\label{ux4f7fux7528ux5143ux7c7b}}

\hypertarget{type}{%
\subsubsection{type()}\label{type}}

动态语言和静态语言最大的不同，就是函数和类的定义，不是编译时定义的，而是运行时动态创建的。

比方说我们要定义一个\texttt{Hello}的
class，就写一个\texttt{hello.py}模块：

\begin{pythoncode}
class Hello(object):
    def hello(self, name='world'):
        print('Hello, %s.' % name)
\end{pythoncode}

当 Python
解释器载入\texttt{hello}模块时，就会依次执行该模块的所有语句，执行结果就是动态创建出一个\texttt{Hello}的
class 对象，测试如下：

\begin{pythoncode}
>>> from hello import Hello
>>> h = Hello()
>>> h.hello()
Hello, world.
>>> print(type(Hello))
<class 'type'>
>>> print(type(h))
<class 'hello.Hello'>
\end{pythoncode}

\texttt{type()}函数可以查看一个类型或变量的类型，\texttt{Hello}是一个
class，它的类型就是\texttt{type}，而\texttt{h}是一个实例，它的类型就是
class \texttt{Hello}。

我们说 class 的定义是运行时动态创建的，而创建 class
的方法就是使用\texttt{type()}函数。

\texttt{type()}函数既可以返回一个对象的类型，又可以创建出新的类型，比如，我们可以通过\texttt{type()}函数创建出\texttt{Hello}类，而无需通过\texttt{class\ Hello(object)...}的定义：

\begin{pythoncode}
>>> def fn(self, name='world'): # 先定义函数
...     print('Hello, %s.' % name)
...
>>> Hello = type('Hello', (object,), dict(hello=fn)) # 创建Hello class
>>> h = Hello()
>>> h.hello()
Hello, world.
>>> print(type(Hello))
<class 'type'>
>>> print(type(h))
<class '__main__.Hello'>
\end{pythoncode}

要创建一个 class 对象，\texttt{type()}函数依次传入 3 个参数：

\begin{enumerate}
\def\labelenumi{\arabic{enumi}.}
\item
  class 的名称；
\item
  继承的父类集合，注意 Python 支持多重继承，如果只有一个父类，别忘了
  tuple 的单元素写法；
\item
  class
  的方法名称与函数绑定，这里我们把函数\texttt{fn}绑定到方法名\texttt{hello}上。
\end{enumerate}

通过\texttt{type()}函数创建的类和直接写 class 是完全一样的，因为 Python
解释器遇到 class 定义时，仅仅是扫描一下 class
定义的语法，然后调用\texttt{type()}函数创建出 class。

正常情况下，我们都用\texttt{class\ Xxx...}来定义类，但是，\texttt{type()}函数也允许我们动态创建出类来，也就是说，动态语言本身支持运行期动态创建类，这和静态语言有非常大的不同，要在静态语言运行期创建类，必须构造源代码字符串再调用编译器，或者借助一些工具生成字节码实现，本质上都是动态编译，会非常复杂。

\hypertarget{metaclass}{%
\subsubsection{metaclass}\label{metaclass}}

除了使用\texttt{type()}动态创建类以外，要控制类的创建行为，还可以使用
metaclass。

metaclass，直译为元类，简单的解释就是：

当我们定义了类以后，就可以根据这个类创建出实例，所以：先定义类，然后创建实例。

但是如果我们想创建出类呢？那就必须根据 metaclass 创建出类，所以：先定义
metaclass，然后创建类。

连接起来就是：先定义 metaclass，就可以创建类，最后创建实例。

所以，metaclass 允许你创建类或者修改类。换句话说，你可以把类看成是
metaclass 创建出来的 ``实例''。

metaclass 是 Python
面向对象里最难理解，也是最难使用的魔术代码。正常情况下，你不会碰到需要使用
metaclass 的情况，所以，以下内容看不懂也没关系，因为基本上你不会用到。

我们先看一个简单的例子，这个 metaclass 可以给我们自定义的 MyList
增加一个\texttt{add}方法：

定义\texttt{ListMetaclass}，按照默认习惯，metaclass 的类名总是以
Metaclass 结尾，以便清楚地表示这是一个 metaclass：

\begin{pythoncode}
class ListMetaclass(type):
    def __new__(cls, name, bases, attrs):
        attrs['add'] = lambda self, value: self.append(value)
        return type.__new__(cls, name, bases, attrs)
\end{pythoncode}

有了 ListMetaclass，我们在定义类的时候还要指示使用 ListMetaclass
来定制类，传入关键字参数\texttt{metaclass}：

\begin{pythoncode}
class MyList(list, metaclass=ListMetaclass):
    pass
\end{pythoncode}

当我们传入关键字参数\texttt{metaclass}时，魔术就生效了，它指示 Python
解释器在创建\texttt{MyList}时，要通过\texttt{ListMetaclass.\_\_new\_\_()}来创建，在此，我们可以修改类的定义，比如，加上新的方法，然后，返回修改后的定义。

\texttt{\_\_new\_\_()}方法接收到的参数依次是：

\begin{enumerate}
\def\labelenumi{\arabic{enumi}.}
\item
  当前准备创建的类的对象；
\item
  类的名字；
\item
  类继承的父类集合；
\item
  类的方法集合。
\end{enumerate}

测试一下\texttt{MyList}是否可以调用\texttt{add()}方法：

\begin{pythoncode}
>>> L = MyList()
>>> L.add(1)
>> L
[1]
\end{pythoncode}

而普通的\texttt{list}没有\texttt{add()}方法：

\begin{pythoncode}
>>> L2 = list()
>>> L2.add(1)
Traceback (most recent call last):
  File "<stdin>", line 1, in <module>
AttributeError: 'list' object has no attribute 'add'
\end{pythoncode}

动态修改有什么意义？直接在\texttt{MyList}定义中写上\texttt{add()}方法不是更简单吗？正常情况下，确实应该直接写，通过
metaclass 修改纯属变态。

但是，总会遇到需要通过 metaclass 修改类定义的。ORM 就是一个典型的例子。

ORM 全称 ``Object Relational Mapping''，即对象 -
关系映射，就是把关系数据库的一行映射为一个对象，也就是一个类对应一个表，这样，写代码更简单，不用直接操作
SQL 语句。

要编写一个 ORM
框架，所有的类都只能动态定义，因为只有使用者才能根据表的结构定义出对应的类来。

让我们来尝试编写一个 ORM 框架。

编写底层模块的第一步，就是先把调用接口写出来。比如，使用者如果使用这个
ORM
框架，想定义一个\texttt{User}类来操作对应的数据库表\texttt{User}，我们期待他写出这样的代码：

\begin{pythoncode}
class User(Model):
    
    id = IntegerField('id')
    name = StringField('username')
    email = StringField('email')
    password = StringField('password')
u = User(id=12345, name='Michael', email='test@orm.org', password='my-pwd')

u.save()
\end{pythoncode}

其中，父类\texttt{Model}和属性类型\texttt{StringField}、\texttt{IntegerField}是由
ORM 框架提供的，剩下的魔术方法比如\texttt{save()}全部由 metaclass
自动完成。虽然 metaclass 的编写会比较复杂，但 ORM
的使用者用起来却异常简单。

现在，我们就按上面的接口来实现该 ORM。

首先来定义\texttt{Field}类，它负责保存数据库表的字段名和字段类型：

\begin{pythoncode}
class Field(object):

    def __init__(self, name, column_type):
        self.name = name
        self.column_type = column_type

    def __str__(self):
        return '<%s:%s>' % (self.__class__.__name__, self.name)
\end{pythoncode}

在\texttt{Field}的基础上，进一步定义各种类型的\texttt{Field}，比如\texttt{StringField}，\texttt{IntegerField}等等：

\begin{pythoncode}
class StringField(Field):

    def __init__(self, name):
        super(StringField, self).__init__(name, 'varchar(100)')

class IntegerField(Field):

    def __init__(self, name):
        super(IntegerField, self).__init__(name, 'bigint')
\end{pythoncode}

下一步，就是编写最复杂的\texttt{ModelMetaclass}了：

\begin{pythoncode}
class ModelMetaclass(type):

    def __new__(cls, name, bases, attrs):
        if name=='Model':
            return type.__new__(cls, name, bases, attrs)
        print('Found model: %s' % name)
        mappings = dict()
        for k, v in attrs.items():
            if isinstance(v, Field):
                print('Found mapping: %s ==> %s' % (k, v))
                mappings[k] = v
        for k in mappings.keys():
            attrs.pop(k)
        attrs['__mappings__'] = mappings 
        attrs['__table__'] = name 
        return type.__new__(cls, name, bases, attrs)
\end{pythoncode}

以及基类\texttt{Model}：

\begin{pythoncode}
class Model(dict, metaclass=ModelMetaclass):

    def __init__(self, **kw):
        super(Model, self).__init__(**kw)

    def __getattr__(self, key):
        try:
            return self[key]
        except KeyError:
            raise AttributeError(r"'Model' object has no attribute '%s'" % key)

    def __setattr__(self, key, value):
        self[key] = value

    def save(self):
        fields = []
        params = []
        args = []
        for k, v in self.__mappings__.items():
            fields.append(v.name)
            params.append('?')
            args.append(getattr(self, k, None))
        sql = 'insert into %s (%s) values (%s)' % (self.__table__, ','.join(fields), ','.join(params))
        print('SQL: %s' % sql)
        print('ARGS: %s' % str(args))
\end{pythoncode}

当用户定义一个\texttt{class\ User(Model)}时，Python
解释器首先在当前类\texttt{User}的定义中查找\texttt{metaclass}，如果没有找到，就继续在父类\texttt{Model}中查找\texttt{metaclass}，找到了，就使用\texttt{Model}中定义的\texttt{metaclass}的\texttt{ModelMetaclass}来创建\texttt{User}类，也就是说，metaclass
可以隐式地继承到子类，但子类自己却感觉不到。

在\texttt{ModelMetaclass}中，一共做了几件事情：

\begin{enumerate}
\def\labelenumi{\arabic{enumi}.}
\item
  排除掉对\texttt{Model}类的修改；
\item
  在当前类（比如\texttt{User}）中查找定义的类的所有属性，如果找到一个
  Field 属性，就把它保存到一个\texttt{\_\_mappings\_\_}的 dict
  中，同时从类属性中删除该 Field
  属性，否则，容易造成运行时错误（实例的属性会遮盖类的同名属性）；
\item
  把表名保存到\texttt{\_\_table\_\_}中，这里简化为表名默认为类名。
\end{enumerate}

在\texttt{Model}类中，就可以定义各种操作数据库的方法，比如\texttt{save()}，\texttt{delete()}，\texttt{find()}，\texttt{update}等等。

我们实现了\texttt{save()}方法，把一个实例保存到数据库中。因为有表名，属性到字段的映射和属性值的集合，就可以构造出\texttt{INSERT}语句。

编写代码试试：

\begin{pythoncode}
u = User(id=12345, name='Michael', email='test@orm.org', password='my-pwd')
u.save()
\end{pythoncode}

输出如下：

\begin{pythoncode}
Found model: User
Found mapping: email ==> <StringField:email>
Found mapping: password ==> <StringField:password>
Found mapping: id ==> <IntegerField:uid>
Found mapping: name ==> <StringField:username>
SQL: insert into User (password,email,username,id) values (?,?,?,?)
ARGS: ['my-pwd', 'test@orm.org', 'Michael', 12345]
\end{pythoncode}

可以看到，\texttt{save()}方法已经打印出了可执行的 SQL
语句，以及参数列表，只需要真正连接到数据库，执行该 SQL
语句，就可以完成真正的功能。

不到 100 行代码，我们就通过 metaclass 实现了一个精简的 ORM
框架，是不是非常简单？

 
 \begin{figure}[htp]
	\centering
	\includegraphics[width=0.6\linewidth]{fig/1183103158580480l.png}
\end{figure}


\hypertarget{ux5c0fux7ed3}{%
\subsubsection{小结}\label{ux5c0fux7ed3}}

metaclass 是 Python
中非常具有魔术性的对象，它可以改变类创建时的行为。这种强大的功能使用起来务必小心。

\hypertarget{ux53c2ux8003ux6e90ux7801}{%
\subsubsection{参考源码}\label{ux53c2ux8003ux6e90ux7801}}

\href{https://github.com/michaelliao/learn-python3/blob/master/samples/oop_advance/create_class_on_the_fly.py}{create\_class\_on\_the\_fly.py}

\href{https://github.com/michaelliao/learn-python3/blob/master/samples/oop_advance/use_metaclass.py}{use\_metaclass.py}

\href{https://github.com/michaelliao/learn-python3/blob/master/samples/oop_advance/orm.py}{orm.py}

